\documentclass{beamer}
\def\presentationtype{0}
\def\corsodilaurea{Ingegneria Informatica e dell'Automazione Industriale}
\def\insegnamento{Ottimizzazione}
\def\nolezione{8}
\def\titolo{Il metodo del simplesso -- Il ``tableau'' e l'operazione ``pivot''}
\def\noattivita{1}

\def\argomenti{Programmazione Lineare}
\def\obiettivilezione{Definire il ``tableau'' e l'operazione ``pivot'' per poter sviluppare metodo del simplesso primale}
\def\nucleotematico{Richiami}
\def\contenuto{Programmazione lineare}


\usepackage[italian]{babel}
%
% Hyperref info for PDF
%
\usepackage{hyperref}
\hypersetup{%
%pdfpagemode=FullScreen,
pdfpagemode=FullScreen,
pdfstartview=Fit,
pdfsubject={\argomenti},
pdfkeywords={\insegnamento,\titolo}
}

\definecolor{ecampus@engred}{RGB}{227,28,23}
\setbeamercolor{normal text}{fg=black,bg=}
\setbeamercolor{alerted text}{fg=red}
\setbeamercolor{example text}{fg=green!50!black}
\setbeamercolor{structure}{fg=ecampus@engred,bg=}
\usetheme{default}
\useinnertheme{rounded}%{circles}
\setbeamertemplate{blocks}[default]
\setbeamercolor{block title}{bg=}
\setbeamercolor{block body}{bg=}

\ifnum\presentationtype>0 {\setbeamertemplate{enumerate item}{Es. \nolezione.\presentationtype.\arabic{enumi})}} \fi

%\setbeamertemplate{frametitle continuation}[from second][(cont'd)]
\usefonttheme[stillsansseriftext,stillsansserifsmall]{serif}
\setbeamertemplate{headline}{%
\leavevmode%
   \hspace*{4.75cm}
   \scalebox{0.90}{%
    \begin{minipage}{\linewidth}%
      \tiny
 	\begin{tabular}{ll}
	  \\
% 	  \textcolor{gray}{Classe}					& \corsodilaurea\\
%	  \textcolor{gray}{Argomento:}					& \insegnamento\\
%	  \textcolor{gray}{Lezione n}:							& \nolezione \\
	  \textcolor{gray}{Titolo:}			& \titolo \\
	  \textcolor{gray}{Attivit\`a}:		& 	  \ifnum\presentationtype=0 {Lezione \nolezione} \else {Sessione di studio \nolezione.\presentationtype} \fi
	\end{tabular}%
    \end{minipage}%
   }%
}
\setbeamertemplate{footline}[text line]{%
  \begin{minipage}{\linewidth}%
    \hfill\insertpagenumber%
%    
    \begin{center}%
      \rule{\linewidth}{0.4pt}
    \end{center}
    \vspace*{-0.65cm}
    \begin{center}%
      \scalebox{0.55}{%
	\begin{minipage}{\linewidth}%
	  \begin{center}%
	    {\tiny
	    \copyright \ 2022 Gionata Massi - \color{blue}{\href{mailto:gionata.massi@savoiabenincasa.it}{gionata.massi@savoiabenincasa.it}}\\
	    \phantom{XXX}}
	  \end{center}
	\end{minipage}%
      }%
    \end{center}%
%
  \end{minipage}%
  }
\setbeamertemplate{navigation symbols}{}

%% Save up on ink for the 4-up handouts
\mode<handout>{%
  \pgfpagesuselayout{4 on 1}[a4paper, landscape, border shrink=10mm]
  \pgfpageslogicalpageoptions{1}{border code=\pgfstroke}
  \pgfpageslogicalpageoptions{2}{border code=\pgfstroke}
  \pgfpageslogicalpageoptions{3}{border code=\pgfstroke}
  \pgfpageslogicalpageoptions{4}{border code=\pgfstroke}
}

\mode<presentation>{\AtBeginSection{%
  \begin{frame}
    \frametitle{Piano della presentazione}
    \tableofcontents[currentsection]
  \end{frame}}
}

\usepackage{microtype}
\usepackage[utf8]{inputenc}
\usepackage[T1]{fontenc}
%\usepackage[osfss]{libertine}
\usepackage{palatino}
\usepackage[scaled=.77]{beramono}
\usepackage{booktabs}
% \usepackage{attachfile}

% Toglie l'enorme spazione prima di description
%\setbeamersize{description width=0pt}

\usepackage{tikz}
\usetikzlibrary{decorations,arrows,shapes,backgrounds,matrix,positioning}
\usepackage{multirow}
\usepackage{verbatim}
\usepackage{eurosym}
\usepackage{pgfpages}

\usepackage{zref-savepos}
\newcounter{restofframe}
\newsavebox{\restofframebox}
\newlength{\mylowermargin}
\setlength{\mylowermargin}{2pt}

\newenvironment{restofframe}{%
    \par%\centering
    \stepcounter{restofframe}%
    \zsavepos{restofframe-\arabic{restofframe}-begin}%
    \begin{lrbox}{\restofframebox}%
}{%
    \end{lrbox}%
    \setkeys{Gin}{keepaspectratio}%
    \raisebox{\dimexpr-\height+\ht\strutbox\relax}[0pt][0pt]{%
    \resizebox*{!}{\dimexpr\zposy{restofframe-\arabic{restofframe}-begin}sp-\zposy{restofframe-\arabic{restofframe}-end}sp-\mylowermargin\relax}%
        {\usebox{\restofframebox}}%
    }%
    \vskip0pt plus 1filll\relax
    \mbox{\zsavepos{restofframe-\arabic{restofframe}-end}}%
    \par
}

\title{\titolo}
\author{Gionata Massi}
\date{\today}

\setbeamercolor*{frametitle}{fg=black}
\setbeamercolor*{title}{fg=black}
\setbeamerfont{frametitle}{shape=\scshape,family=\rmfamily,size=\large,series=\bfseries}

\setbeamersize{}
\makeatletter
\newcommand\thefontsize[1]{{#1 The current font size is: \f@size pt\par}}
\makeatother

\tikzstyle{nicebox}=[draw=gray!100, fill=blue!10, very thick,
rounded corners, rectangle, inner sep=4pt, inner ysep=16pt]
\tikzstyle{niceboxtitle}=[draw=gray!100, fill=white, text=black,
rounded corners, very thick, rectangle]
\newcommand\nicebox[2]{
{\centering
\begin{tikzpicture}
\node [nicebox](box){
\begin{minipage}{0.95\textwidth}\centering
\begin{minipage}{0.95\textwidth}
#2
\end{minipage}\end{minipage}};
\node[niceboxtitle, right=10pt] at (box.north west)
{\small\textbf{#1}};
\end{tikzpicture}\par}
}

\tikzstyle{modelbox}=[draw=structure!100, fill=white!50, very thick,
rounded corners, rectangle, inner sep=4pt, inner ysep=16pt, text=blue!50!black]
\tikzstyle{modelboxtitle}=[draw=structure!100, fill=white!50, text=blue!50!black,
rounded corners, very thick, rectangle]
\newcommand\modelbox[2]{
{\centering
\begin{tikzpicture}
\node [modelbox](box){
\begin{minipage}{0.95\textwidth}\centering
\begin{minipage}{0.95\textwidth}
#2
\end{minipage}\end{minipage}};
\node[modelboxtitle, right=10pt] at (box.north west)
{\small\textbf{\mbox{#1}}};
\end{tikzpicture}\par}
}

% % %
% Tipi di sessioni di studio
% % %
\def\domande{Domande}
\def\approfondimenti{Approfondimenti}
\def\esercizi{Esercizi}
\def\esempi{Esempi svolti}
\def\relazioni{Relazioni}
\def\attivitapratica{Attivit\`a pratica}
% % %

\newcommand{\generatitolo}{ %
\begin{frame}[plain]
\pdfbookmark{\ifnum\presentationtype=0 {Lezione \nolezione} \else {Sessione di studio \nolezione.\presentationtype} \fi}{\ifnum\presentationtype=0 {Lezione \nolezione} \else {Sessione di studio \nolezione.\presentationtype} \fi}
\noindent\hspace*{3.715cm}
   \scalebox{0.90}{%
    \begin{minipage}{\linewidth}%
      \tiny
 	\begin{tabular}{ll}
	  \\
% 	  \textcolor{gray}{Classe}					& \corsodilaurea\\
%	  \textcolor{gray}{Argomento:}					& \insegnamento\\
%	  \textcolor{gray}{Lezione n}:							& \nolezione \\
	  \textcolor{gray}{Titolo:}			& \titolo \\
	  \textcolor{gray}{Attivit\`a}:		& 	  \ifnum\presentationtype=0 {Lezione \nolezione} \else {Sessione di studio \nolezione.\presentationtype} \fi
	\end{tabular}%
    \end{minipage}%
   }%
   
  \begin{center}
    \large{\textbf{\insegnamento}}
  \end{center}

  \vspace*{0.5cm}
  
  \begin{center}
    \huge{\titolo}
  \end{center}
  
  \begin{center}
	  \ifnum\presentationtype=0
	    {Lezione n$^{\circ}$ \nolezione}
	  \else
	    {Sessione di studio n$^{\circ}$ \nolezione .\presentationtype}
	  \fi
  \end{center}
  
  \vspace*{0.5cm}

  \begin{center}
    \large{\textsl{Gionata Massi}}\\
    \scriptsize\color{blue}{<\href{mailto:gionata.massi@savoiabenincasa.it}{gionata.massi@savoiabenincasa.it}>}
  \end{center}
  
  \newcounter{pageleft}
  \zsaveposy{pageleft}
  \vspace*{\dimexpr\zposy{pageleft}sp-1.67cm}

  \begin{minipage}{\linewidth}%
   \hfill\phantom{\insertpagenumber}%

  \begin{center}%
      \rule{\linewidth}{0.4pt}
    \end{center}
    
    \vspace*{-0.95cm}
    
    \begin{center}%
      \scalebox{0.55}{%
	\begin{minipage}{\linewidth}%
	  %\begin{center}%
	  %  \tiny
	  %  \copyright \ 2022 Gionata Massi - \color{blue}{\href{mailto:gionata.massi@savoiabenincasa.it}{gionata.massi@savoiabenincasa.it}}\\
	  %  \color{blue}{\href{mailto:gionata.massi@savoiabenincasa.it}{gionata.massi@savoiabenincasa.it}}\\
	  %  \phantom{XXX}
	  %\end{center}
	\end{minipage}%
      }%
    \end{center}%
     \phantom{XXX}
%
  \end{minipage}%

\end{frame}
}

\renewcommand{\vec}{\mathbf}
\newcommand{\matr}{\mathbf}

\uselanguage{italian}
\languagepath{italian}
\deftranslation[to=italian]{Theorem}{Teorema}
\deftranslation[to=italian]{theorem}{teorema}
\deftranslation[to=italian]{Definition}{Definizione}
\deftranslation[to=italian]{definition}{definizione}
\deftranslation[to=italian]{Corollary}{Corollario}
\deftranslation[to=italian]{corollary}{corollario}

\let\definition\relax
\let\enddefinition\relax

%\theoremstyle{example}
\newtheorem{definition}[theorem]{\translate{Definition}}

\def\fieldN{\mathbb{N}}
\def\fieldR{\mathbb{R}}
\def\fieldZ{\mathbb{Z}}

\def\matrA{\matr{A}}
\def\vecB{\vec{b}}
\def\vecC{\vec{c}}
\def\vecX{\vec{x}}

\usepackage{cancel}
\usepackage{bookmark}

% tkiz ball item
\newcommand*\circled[1]{\raisebox{2.5pt}{\tikz[baseline=(char.base)]{
            \node[circle,ball color=purple, shade, 
 color=white,inner sep=1.2pt] (char) {\tiny #1};}}}

% tkiz rounded item
\newcommand*\rounded[1]{\tikz[baseline=(char.base)]{
            \node[draw=none,ball color=purple, shade, 
 color=white, rounded corners=3.5pt, inner sep=2.5pt] (char) {\scriptsize #1};}}

 \usepackage{etoolbox}
\makeatletter
\apptocmd{\beamer@@frametitle}{\write\@auxout{\string\@writefile{frm}{\string\frametitleentry{\the\c@framenumber}{#1}{#2}}}}{}{}
\newcommand*{\frametitleentry}[3]{\@namedef{frametitleshort#1}{#2}\@namedef{frametitle#1}{#3}}
\AtEndDocument{\if@filesw\newwrite\tf@frm\immediate\openout\tf@frm\jobname.frm\relax\fi}
\@input{\jobname.frm}
\makeatother
 
%\usebackgroundtemplate{\includegraphics[width=\paperwidth]{../img/sfondo_ecampus.png}}

\usepackage{listings}
% Extension for listings package
\lstdefinelanguage{ampl}{
  % See AMPL book (1993 edition), Appendix A
  morekeywords={%
    % Table A.1
    % Excluding symbols and "not X", where X (and not) are keywords
    if, then, else, or,  exists, forall,  and, in, within, 
    not, union, diff, symdiff, inter, cross, setof, by, 
    % Skipping: less,  sum, prod, min, max, div, mod,
    % Skipping built-in functions (Tables A-2 and A-3)
    % Section A.5 -- Declarations of model entities (subject to ->
    % subject, to)
    set, param, var, arc, minimize, maximize, subject, to, node,
    % Section A.6 -- Set declarations (not repeating set or within)
    dimen, default, card, ordered, by, reversed, circular,
    % Section A.7 -- Parameter declarations (not repeating in, default)
    param, binary, integer, symbolic, check, Infinity
    % Section A.8 -- Variable declarations
    var,  coeff, cover, obj,
    % Section A.13 -- Command language
    % Skipping most commands!
    include, model, data, let, objective, drop, restore, solve,
    solution, quit, end, option, reset, update, commands,
    % A.14 -- Imported functions
    function, pipe, 
    % A.17 synonyms for previously-defined keywords
    coef, difference, dimen, dimension, intersect, intersection, s.t., 
    % 
    % must be defined somewhere ...
    for, while, repeat, break,
    % NEW language features (http://www.ampl.com/NEW/newlang.html)
    % Grouped by page
    purge, problem, redeclare, xref,
    complements, 
    table, IN, OUT, INOUT, read, write, 
    expand,
    for, commands, repeat, until, while, break, continue,
  },
  sensitive=true,
  comment=[l]{\#},
  morecomment=[s]{/*}{*/},
  string=[b]",
  morestring=[b]',
}[keywords,comments,strings]

\definecolor{comment_c}{RGB}{60,128,49}
\definecolor{rowname_c}{RGB}{128,60,49}
\definecolor{reserved_c}{RGB}{49,60,128}

\lstset{
	language=ampl,%
	commentstyle=\color{comment_c},%
	backgroundcolor=\color{blue!15},%
    basicstyle=\small\ttfamily,%
    %emph={s.t.},emphstyle={\bfseries},%
    otherkeywords={s.t.},%
    numbers=left, numberstyle=\tiny, stepnumber=1, numbersep=5pt,%
}

\definecolor{col_set}{RGB}{128,0,0}
\definecolor{col_par}{RGB}{91,0,91}
\definecolor{col_idx}{RGB}{0,91,91}
\definecolor{col_var}{RGB}{0,0,128}

\def\insieme{%1}{\rm \textcolor{col_set}{%1}}}


\begin{document}

\generatitolo

\section{Sistemi di equazioni lineari e matrici}

\begin{frame}[allowframebreaks]{Forma generale di un sistema lineare}
  
  \modelbox{Sistema lineare di $m$ equazioni in $n$ incognite}{%
  \[%
% \left\{\begin{array}{lclccclcr}
%   a_{11} x_1 &+& a_{12} x_2 &+& \cdots &+& a_{1n} x_n &=& b_1\\
%   a_{21} x_1 &+& a_{22} x_2 &+& \cdots &+& a_{2n} x_n &=& b_2\\
%              & &            & & \ddots & &            &\vdots& \vdots\\
%   a_{m1} x_1 &+& a_{m2} x_2 &+& \cdots &+& a_{mn} x_n &=& b_m
% \end{array}\right.%
\left\{
\begin{matrix} 
a_{11}x_1 + a_{12}x_2 +\cdots + a_{1n}x_n = b_1\\
a_{21}x_1 + a_{22}x_2 +\cdots + a_{2n}x_n = b_2\\
\vdots\\
a_{m1}x_1 + a_{m2}x_2 +\cdots + a_{mn}x_n = b_m\\
\end{matrix}
\right.
\]}

  \framebreak
  
  \begin{itemize}
  \item  Il numero $n$ delle incognite \`e detto ordine del sistema

  \item  I numeri $a_{11},\ a_{12}, \ldots,\ a_{mn}$  sono i coefficienti del sistema

  \item  I numeri $b_1,\ b_2, \ldots,\ b_m$ sono i termini noti del sistema
  
  \item  Noti $m$ ed $n$ \`e possibile scrivere un sistema di equazioni lineari usando
  la notazione matriciale
  \end{itemize}
  
  \framebreak
  
  I coefficienti del sistema possono essere raccolti in una tabella di numeri detta \emph{matrice dei coefficienti}

  \modelbox{Matrice dei coefficienti}{%
  \[
  \matr{A} = 
  \begin{bmatrix}\begin{array}{cccc}
  a_{11} & a_{12} & \cdots & a_{1n}\\
  a_{21} & a_{22} & \cdots & a_{2n}\\
         &        & \ddots \\
  a_{m1} & a_{m2} & \cdots & a_{mn}
  \end{array}\end{bmatrix}
  \]
  }

  \framebreak
  
   I termini noti del sistema possono essere raccolti nel vettore colonna detto \emph{vettore dei termini noti}

  \modelbox{Vettore dei termini noti}{%
  \[
  \vec{b} = 
  \begin{bmatrix}\begin{array}{c}
  b_{1}\\
  b_{2}\\
  \vdots \\
  b_{m}
  \end{array}\end{bmatrix}
  \]
  }
  
  \framebreak
  
  Le incognite possono essere rappresentate dal \emph{vettore delle incognite}

  \modelbox{Vettore delle incognite}{%
  \[
  \vec{x} = 
  \begin{bmatrix}\begin{array}{c}
  x_{1}\\
  x_{2}\\
  \vdots \\
  x_{n}
  \end{array}\end{bmatrix}
  \]
  }

  \small{L'intero sistema lineare  di $m$ equazioni in $n$ incognite \`e rappresentato dalla \emph{matrice completa} (o \emph{matrice aumentata}) del sistema}

  \modelbox{Matrice completa}{%
  \[
  [\matr{A}|\vec{b}] = 
  \begin{bmatrix}
  \left.\begin{array}{cccc}
  a_{11} & a_{12} & \cdots & a_{1n}\\
  a_{21} & a_{22} & \cdots & a_{2n}\\
         &        & \ddots \\
  a_{m1} & a_{m2} & \cdots & a_{mn}
  \end{array}\right|
  \begin{array}{c}
  b_{1}\\
  b_{2}\\
  \vdots\\
  b_{m}
  \end{array}
  \end{bmatrix}
  \]
  }
\end{frame}

\section{Prodotto di una matrice per un vettore}

\begin{frame}{Matrici}
\footnotesize{
  \begin{itemize}
  \item In generale una matrice con $m$ righe e $n$ colonne \`e
  una tabella rettangolare di numeri con $m$ righe e $n$
  colonne
  
  \item L'insieme delle matrici con $m$ righe e $n$ colonne ($m$x$n$) a coefficienti
  reali \`e indicato con $M_{m,n}(\mathbb{R})$

  \item Sia $\matr{A} \in M_{m,n}(\mathbb{R})$ una matrice $m$x$n$
  \begin{description}
   \item[$a_{ij}$] \`e il coefficiente di posto $(i, j)$ della matrice $\matr{A}$
   \item[$i$]  \`e l'indice di riga
   \item[$j$] \`e l'indice di colonna
   \item[$A_i$] \`e la riga $i$-esima; $A_i=\left[\begin{array}{cccc}a_{i1}&a_{i2}&\cdots&a_{in}\end{array}\right]$
   \item[$A^j$] \`e la colonna $j$-esima; $A^j=\left[\begin{array}{c}a_{1j}\\a_{2j}\\\vdots\\a_{mj}\end{array}\right]$
  \end{description}
  \end{itemize}
}
\end{frame}

\begin{frame}[allowframebreaks]{Prodotto di una matrice per un vettore}
Si definisce il prodotto di una matrice $\matr{A} \in M_{m,n}(\mathbb{R})$
per un vettore $x \in \mathbb{R}^n$ come:

\modelbox{Prodotto di una matrice per un vettore}{%
  \small\[
  \matr{A} \vec{x} = 
  \begin{bmatrix}\begin{array}{cccc}
  a_{11} & a_{12} & \cdots & a_{1n}\\
  a_{21} & a_{22} & \cdots & a_{2n}\\
         &        & \ddots \\
  a_{m1} & a_{m2} & \cdots & a_{mn}
  \end{array}\end{bmatrix}
  \begin{bmatrix}\begin{array}{c}
  x_{1}\\
  x_{2}\\
  \vdots \\
  x_{n}
  \end{array}\end{bmatrix} =
  x_{1} A^1 + x_{2} A^2 + \cdots + x_{n} A^n 
  \]
}

\framebreak

 che può essere riscritto come il vettore 

\modelbox{Prodotto di una matrice per un vettore}{%
  \small\[
  \matr{A} \vec{x} = 
  \begin{bmatrix}\begin{array}{cccc}
  a_{11} x_{1} +& a_{12} x_{2} +& \cdots +& a_{1n} x_{n}\\
  a_{21} x_{1} +& a_{22} x_{1} +& \cdots +& a_{2n} x_{n}\\
         &        & \ddots \\
  a_{m1} x_{1} +& a_{m2} x_{2} +& \cdots +& a_{mn} x_{n}
  \end{array}\end{bmatrix}
  \]
}
\end{frame}

\begin{frame}{Forma matriciale di un sistema di equazioni lineari}

\scriptsize{
 Dalla precedente definizione di prodotto di una matrice per un vettore abbiamo che il sistema di equazioni lineari

\[
\left\{\begin{array}{lclccclcc}
  a_{11} x_1 &+& a_{12} x_2 &+& \cdots &+& a_{1n} x_n &=& b_1\\
  a_{21} x_1 &+& a_{22} x_2 &+& \cdots &+& a_{2n} x_n &=& b_2\\
             & &            & & \ddots & &            &\vdots& \vdots\\
  a_{m1} x_1 &+& a_{m2} x_2 &+& \cdots &+& a_{mn} x_n &=& b_m
\end{array}\right.
\]

può essere espresso come

\[
\begin{bmatrix}\begin{array}{cccc}
  a_{11} x_{1} +& a_{12} x_{2} +& \cdots +& a_{1n} x_{n}\\
  a_{21} x_{1} +& a_{22} x_{1} +& \cdots +& a_{2n} x_{n}\\
         &        & \ddots \\
  a_{m1} x_{1} +& a_{m2} x_{2} +& \cdots +& a_{mn} x_{m}
\end{array}\end{bmatrix} 
=
  \begin{bmatrix}\begin{array}{c}
  b_{1}\\
  b_{2}\\
  \vdots \\
  b_{m}
  \end{array}\end{bmatrix}
\]

o anche: \begin{center}{$\matr{A} \vec{x} = \vec{b}$}\end{center}
}
\end{frame}

\section{Definizioni e teoremi di Algebra Lineare}

\begin{frame}[allowframebreaks]{Definizioni}

\begin{definition}[soluzione]
   Una soluzione del sistema di equazioni lineari \`e una
   $n$-upla $(v_1, v_2, \ldots, v_n)$ di numeri che sostituiti
   ordinatamente alle incognite $(x_1, x_2, \ldots, x_n)$
   soddisfano  tutte le equazioni del sistema.
\end{definition}

\begin{definition}[combinazione lineare]
  Una combinazione lineare dei vettori $v_1, v_2, \ldots, v_k \in V$
  \`e un vettore $w = \alpha_1 v_1+ \alpha_2 v_2 +\ldots+ \alpha_k v_k\quad\alpha_j \in \mathbb{R} \text{ per } j = 1, 2, \ldots, k$
\end{definition}

\framebreak

\begin{definition}[dipendenza lineare]
  Un insieme di vettori $v_1, v_2, \cdots, v_k \in V$ si  dicono
  linearmente dipendenti se esistono 
  $\alpha_1\in \mathbb{R}, \alpha_2\in \mathbb{R},\ldots, \alpha_k \in \mathbb{R}$
  non tutti nulli tali che
  $\alpha_1 v_1+ \alpha_2 v_2 +\cdots+ \alpha_k v_k = 0 $
\end{definition}

\begin{definition}[lineare indipendenza]
  Un insieme di vettori $v_1, v_2, \cdots, v_k$ si dicono linearmente indipendenti se 
  $\alpha_1 v_1+ \alpha_2 v_2 +\cdots+ \alpha_k v_k = 0$ implica
  $\alpha_1 = \alpha_2 = \cdots = \alpha_k = 0$
\end{definition}

\begin{definition}[base]
  Una base per un insieme di vettori \`e un sottoinsieme di vettori linearmente indipendenti tali che ogni altro vettore dell'insieme può essere espresso come una loro combinazione lineare.
\end{definition}

\begin{definition}[rango]
  Il rango di un insieme di vettori \`e il più grande numero di vettori linearmente indipendenti che si possono scegliere nell'insieme.

  Il rango di riga di una matrice quadrata \`e il rango dell'insieme dei suoi vettori riga. Il rango di colonna di una matrice quadrata \`e il rango dell'insieme dei suoi vettori colonna. 

  Il rango di riga e il rango di colonna di una matrice sono uguali tra loro e sono uguali al numero di pivot (calcolabili con una qualsiasi eliminazione di Gauss) 
\end{definition}

\begin{definition}[singolarit\`a]
  Una matrice \`e detta non singolare se il rango \`e uguale sia al suo numero di righe che al suo numero di colonne, cio\`e possiede tutti i pivot. E detta singolare altrimenti.
\end{definition}

\begin{definition}[sottomatrice]
   Una sottomatrice di una matrice
   $\matr{A} \in M_{m,n}(\mathbb{R})$
   \`e una matrice
   $\matr{A} \in M_{r,s}(\mathbb{R})$
   ottenuta da $\matr{A}$
   rimuovendo $m - r$ righe e $n - s$ colonne
\end{definition}
   
\begin{definition}[minore]
  Un minore \`e una sottomatrice quadrata, cio\`e con $r = s$.
  Il numero $r$ \`e definito ordine del minore
\end{definition}
   
\begin{definition}[minore complementare]  
  Un minore complementare  di una matrice
  $\matr{A} \in M_{m,n}(\mathbb{R})$  \`e un suo minore ottenuto togliendo
  una sola riga e una sola colonna.
  Il minore ottenuto togliendo l'$i$-esima riga e la $j$-esima colonna si indica con
  $A(i,j)$
\end{definition}
   
\begin{definition}[rango di una matrice rettangolare]
  Se una matrice non \`e quadrata il rango \`e il
  massimo ordine di un minore non singolare
\end{definition}
\end{frame}

\begin{frame}[allowframebreaks]{Teoremi}
\begin{theorem}[Rouché-Capelli]
 dato un sistema di equazioni lineari
 $\matr{A}\vec{x} = \vec{b}$, con
 $\matr{A} \in M_{m,n}(\mathbb{R}),\ \vec{x} \in \mathbb{R}^n$ e $\vec{b} \in \mathbb{R}^m$
  esistono soluzioni per il sistema se
  e solo se il rango della matrice completa $[\matr{A}|\vec{b}]$
  è uguale al rango della matrice $\matr{A}$;
  se esistono soluzioni e il rango di $\matr{A}$
  è uguale a $n$, allora la soluzione è unica;
  se esistono soluzioni e il rango di $\matr{A}$
  è minore di $n$, esistono infinite soluzioni
\end{theorem}
il teorema afferma che le soluzioni formano
  un sottospazio affine di dimensione $n - \operatorname{rango}\matr{A}$
\framebreak

\begin{theorem}[Regola di Cramer]
    dato un sistema di equazioni lineari
    $\matr{A}\vec{x} = \vec{b}$, con
    $\matr{A} \in M_{n,n}(\mathbb{R})$ e invertibile,
    $\vec{x} \in \mathbb{R}^n$  e $\vec{b} \in \mathbb{R}^n$
    è possibile calcolare gli elementi del vettore soluzione come

    \[x_i=\frac{\det(\matr{A}(i))}{\det(\matr{A}	)}\]
 
    dove $\matr{A}(i)$ è la matrice formata sostituendo la $i$-esima colonna di
    $\matr{A}$ con il vettore $\vec{b}$.

\end{theorem}

\framebreak

\begin{exampleblock}{Sistemi di equazioni rettangolari $m < n$}
    Dato un sistema di $m$ equazioni in $n$ incognite ($m$ < $n$)
    
  sia $\operatorname{rango}\matr{A} = \operatorname{rango}[\matr{A}|\vec{b}]$;
  
  sia $\operatorname{rango}\matr{A} = m$ (se non è uguale si eliminino le equazioni linearmente dipendenti);
  
      allora il sistema ha infinite soluzioni in quanto $n - m$ variabili sono arbitrarie.
 \end{exampleblock}
\end{frame}

\end{document}
